\section{Optional application: normal-phase sound damping}
\label{qt:hilker}
\subsection{What is compared with the experiment?}
An experimental application requires more than changing the statistical sign.
We need the particles' scattering cross section, their thermodynamic state and
a model for the measured observable. We use the normal-phase sound measurements
in Fig.~4 of Hilker et al.\ \cite{qtHilker}, with the public data in
Ref.~\cite{qtHilkerData}. We do not model their superfluid second sound or
claim to reproduce the complete experiment.

The chain of calculations is
\[
 \text{scattering model}
 \longrightarrow(\mu,\lambda)_{\mathrm{Galerkin}}
 \longrightarrow\text{linear fluid equations}
 \longrightarrow\text{sound speed and damping}.
\]
In this section only, the standard sound-literature symbols
$\eta=\mu$ and $\kappa=\lambda$ avoid confusing viscosity with chemical
potential. Sound propagation uses a macroscopic closure informed by the
computed transport coefficients. It is distinct from directly evolving the
projected transverse kinetic equation of the preceding section.

\subsection{An experimentally motivated kernel and units}
For low-energy identical-boson scattering with scattering length $a$, the
isotropic total cross section is $8\pi a^2$. In our full-solid-angle convention,
\begin{equation}
 B=2a^2g,\qquad d^2=8a^2.
\end{equation}
The same dimensionless Galerkin solver can therefore be used, with
\begin{equation}
 \eta=\mu_{\mathrm{num}}\frac{\sqrt{mk_{\mathrm B}T}}{8a^2},\qquad
 \kappa=\lambda_{\mathrm{num}}
              \frac{k_{\mathrm B}\sqrt{k_{\mathrm B}T/m}}{8a^2}.
 \label{qt:atomunits}
\end{equation}
This assumes the low-energy, energy-independent cross section remains an
adequate description. Finite-energy scattering or many-body corrections
would modify this model rather than the statistical sign alone.

The selected normal-phase data concern $^{39}$K at approximately $1.3T_c$.
We retain five observations with the published inverse Knudsen number
$I=L/\ell\geq3$, favouring the more hydrodynamic part of the dataset.
The density is reconstructed using the published derived quantity:
\begin{equation}
 n=\frac{I}{8\pi a^2L}.
 \label{qt:densitydata}
\end{equation}
This makes the comparison conditional on that density reconstruction; it is
not an independent determination of density from raw experimental images.
One additional source row has inconsistent length conventions when its
reported damping ratio is checked (the listed $70\,\mu\mathrm m$ corresponds
to about $72\,\mu\mathrm m$ in that ratio). It is retained unchanged in the
data archive and already lies outside the selected hydrodynamic subset.

\subsection{Why the thermodynamic model must also change}
The ideal-gas kinetic pressure alone is not the thermodynamic model used for
this weakly interacting Bose experiment. In the normal Hartree--Fock
approximation, introduce
\begin{equation}
 g_{\mathrm{int}}=\frac{4\pi\hbar^2a}{m},\qquad
 z=\exp\!\left(\frac{\mu_{\mathrm{chem}}-2g_{\mathrm{int}}n}
                         {k_{\mathrm B}T}\right).
\end{equation}
The effective single-particle energy is shifted by $2g_{\mathrm{int}}n$.
The occupation still has the Bose form in terms of the effective fugacity
$z$, while the total pressure and energy density acquire interaction terms:
\begin{equation}
 p_{\mathrm{kin}}=nk_{\mathrm B}T\frac{G_{5/2}(z)}{G_{3/2}(z)},\qquad
 p=p_{\mathrm{kin}}+g_{\mathrm{int}}n^2,\qquad
 \mathcal U=\frac32p_{\mathrm{kin}}+g_{\mathrm{int}}n^2.
\end{equation}
The kinetic transport source still contains $Z_{\mathrm{kin}}$, not
$p/(nk_{\mathrm B}T)$ including the interaction pressure. To see why,
the mean-field acceleration is $\mathbf X=-2g_{\mathrm{int}}\nabla n/m$.
The Euler equation for the total pressure then gives
\[
 D_t\mathbf u-\mathbf X
 =-\frac{\nabla(p_{\mathrm{kin}}+g_{\mathrm{int}}n^2)}{mn}
   +\frac{2g_{\mathrm{int}}}{m}\nabla n
 =-\frac{\nabla p_{\mathrm{kin}}}{mn}.
\]
This cancellation preserves the matched thermal source used in the collision
calculation. Adding an interaction pressure without its corresponding force
would be inconsistent.

Let $p_n=(\partial p/\partial n)_T$ and $p_T=(\partial p/\partial T)_n$.
The derivatives used by the sound calculation are
\begin{align}
 p_n&=k_{\mathrm B}T\frac{G_{3/2}}{G_{1/2}}+2g_{\mathrm{int}}n,\\
 p_T&=nk_{\mathrm B}\left(\frac52\frac{G_{5/2}}{G_{3/2}}
                  -\frac32\frac{G_{3/2}}{G_{1/2}}\right),\\
 C_V&=\frac32p_T,\qquad
 C_P=C_V+\frac{Tp_T^2}{np_n},\qquad
 c_{\mathrm{ad}}^2=\frac{C_P}{C_V}\frac{p_n}{m}.
 \label{qt:thermo}
\end{align}
Here $C_V$ and $C_P$ are heat capacities per unit volume. Their units matter:
they must not be mixed with heat capacities per unit mass or per particle.

\subsection{Deriving the sound eigenproblem}
Linearize continuity, longitudinal momentum and temperature about a uniform
state. With $\rho=mn$ and zero bulk viscosity, they read
\begin{align}
 \partial_t\delta n&=-n\partial_xu,\\
 \rho\partial_tu&=-\partial_x(p_n\delta n+p_T\delta T)
                         +\frac43\eta\partial_x^2u,\\
 C_V\partial_t\delta T&=-Tp_T\partial_xu+\kappa\partial_x^2\delta T.
 \label{qt:soundeq}
\end{align}
The term $4\eta/3$ is the longitudinal projection of the same traceless
viscous stress derived earlier. Thermal conduction couples temperature to
the density and velocity oscillation, so both computed coefficients contribute
to sound damping.

Insert $e^{i(kx-\omega t)}$, take $k=\pi/L$ for the fundamental mode and
introduce
\begin{equation}
 C=\frac\omega{kc_{\mathrm{ad}}},\qquad
 q_\eta=\frac{4\eta k}{3\rho c_{\mathrm{ad}}},\qquad
 q_\kappa=\frac{\kappa k}{C_Vc_{\mathrm{ad}}}.
\end{equation}
Eliminating the three amplitudes gives
\begin{equation}
 C^3+i(q_\eta+q_\kappa)C^2
 -(1+q_\eta q_\kappa)C
 -iq_\kappa\frac{C_V}{C_P}=0.
 \label{qt:cubic}
\end{equation}
We select the damped propagating sound root. The quantities compared with
the published data conventions are
\begin{equation}
 c_I=\frac{|\omega|}{k},\qquad \Gamma=-\operatorname{Im}\omega.
\end{equation}
Thus $c_I$ is not simply $\operatorname{Re}\omega/k$ when damping is
appreciable. The dataset quotes $\Gamma/(2\pi)$ in hertz; the factor $2\pi$
must be retained when converting the eigenvalue. In the weak-damping limit,
\begin{equation}
 \Gamma\simeq\frac{k^2}{2}
 \left[\frac{4\eta}{3\rho}+\kappa\left(\frac1{C_V}-\frac1{C_P}\right)\right].
\end{equation}
The code uses the full cubic rather than imposing this last asymptotic formula.

\subsection{Results and the limits of this comparison}
\begin{figure}[htbp]\centering\qtfig{hilker}
\caption{Five selected normal-phase observations from Hilker et al., plotted
against inverse Knudsen number. Circle and triangle markers distinguish
$L=50$ and $70\,\mu\mathrm m$ boxes. Colors distinguish measurements,
quantum transport and classical transport, the latter two using identical
Hartree--Fock thermodynamics and linear Navier--Stokes--Fourier sound.
Data bars are reported measurement uncertainties. Quantum-model bars are
one-input-at-a-time sensitivity ranges, not confidence intervals;
classical markers show central predictions only. No lines connect states
with different physical inputs. Data: Ref.~\cite{qtHilkerData}, CC BY 4.0.}
\label{qt:hilkerfig}\end{figure}
\begin{table}[htbp]\centering\small
\begin{tabular}{rrrrrrr}\toprule
 & & &\multicolumn{2}{c}{$c_I$ [mm/s]}&\multicolumn{2}{c}{$\Gamma/(2\pi)$ [Hz]}\\
Index&$L$ [$\mu$m]&$a/a_{\mathrm B}$&Data&Model&Data&Model\\\midrule
\qtinput{\qtroot/data/hilker_table.tex}
\bottomrule\end{tabular}
\caption{Central values for the observations in Fig.~\ref{qt:hilkerfig}.
$a_{\mathrm B}$ is the Bohr radius. No parameters were fitted to these
sound measurements.}\label{qt:hilkertable}
\end{table}
For this extension we increase the basis to $N=7$. Across the five states,
the largest change from $N=5$ to $N=7$ is about $0.0351\%$ in viscosity and
$0.1679\%$ in conductivity. Refining the collision quadrature from
$(24,16,20)$ to $(32,20,24)$ changes these coefficients by less than
$0.001\%$. These numerical changes are much smaller than the central-value
differences from experiment: approximately $-8.5\%$ to $+3.5\%$ for $c_I$
and $-20.9\%$ to $+15.5\%$ for damping.

\paragraph{Classical-transport control.}
We repeat the sound prediction with classical viscosity and conductivity
at the same basis size and physical state, leaving the Hartree--Fock
thermodynamics and cross-section convention unchanged. This isolates the
transport correction; it is not a fully classical thermodynamic model.
For either measured observable $y$, define the descriptive normalized RMS
residual
\[
 R_y=\left[\frac{1}{5}\sum_{j=1}^{5}
 \left(\frac{y_j^{\mathrm{model}}-y_j^{\mathrm{data}}}
 {\sigma_{y,j}^{\mathrm{data}}}\right)^2\right]^{1/2}.
\]
The report exporter computes the following values from the stored predictions:
\begin{center}
\begin{tabular}{lrr}\toprule
Model (same HF thermodynamics)&$R_{c_I}$&$R_\Gamma$\\\midrule
\qtinput{\qtroot/data/hilker_metrics_table.tex}
\bottomrule\end{tabular}
\end{center}
The RMS agreement is nearly unchanged. Although the largest relative damping
error decreases from $34.2\%$ for classical transport to $20.9\%$ for quantum
transport, these five observations do not demonstrate an overall improvement
from quantum transport corrections. Neither calculation fits the measured
sound values. The residual normalization uses only the quoted measurement
uncertainty: it does not include propagated input uncertainty or covariance
and is not a formal significance test. The value $R_{c_I}=1.031$ is an aggregate
and does not mean that every point agrees within that many uncertainties.

The comparison therefore provides a useful application with a visible
remaining discrepancy, rather than precision experimental validation of the
collision solver. Possible limitations include the hydrodynamic closure at
finite $L/\ell$, the normal-phase thermodynamic approximation, the assumed
cross section and the reconstruction of state variables. Their individual
contributions have not been separated. Matching a separately published theory
curve is another algebraic/unit check and should not be counted as matching
the measured data.

The repository also evolves an illustrative standing wave initialized with
a $1\%$ density perturbation and zero initial velocity and temperature
perturbation. This illustrates Eq.~\eqref{qt:soundeq}; it does not reproduce
the experimental shaking and preparation protocol. The main result of the
project remains the verified collision-to-transport computation. This final
application shows how those coefficients enter an experimentally measurable
quantity, and where further physical modelling becomes necessary.
